\pdfoutput=1
\documentclass[11pt,a4paper]{article}
\usepackage{abhijit-research}
\renewcommand{\PaperCode}{P08}
\renewcommand{\PaperKind}{RESEARCH PAPER}
\renewcommand{\PaperVersion}{VERSION 1.1}
\renewcommand{\PaperDate}{AUGUST 2026}
\renewcommand{\PaperTitle}{Critical Capacity Completion}
\renewcommand{\PaperSubtitle}{Equivalent Forms of the Riemann Hypothesis}
\renewcommand{\PaperFullTitle}{Critical Capacity Completion and Equivalent Forms of the Riemann Hypothesis}
\renewcommand{\PaperShortTitle}{Critical Capacity Completion}
\renewcommand{\PaperKeywords}{Riemann Hypothesis; Mobius function; Mertens function; capacity energy; operator inequality; critical exponent}
\renewcommand{\PaperClassification}{The finite reconstruction, weighted-energy divergence, and equivalence theorems are asserted with their displayed proofs and scopes. The universal RH-equivalent lower bound remains open; no proof of RH is claimed.}
\renewcommand{\PaperCitation}{Singh, Abhijit. 2026. \textit{Critical Capacity Completion and Equivalent Forms of the Riemann Hypothesis}. Research manuscript, version 1.0.}
\begin{document}
\MakeResearchTitle
\MakeResearchContents
\MakeResearchAbstract{%
A frozen source-derived arithmetic construction is connected to the classical Riemann problem through exact finite and infinite statements. Scalarization of the uncollapsed response algebra fails only through repeated-prime and over-capacity interactions; differentiation gives an exact boundary-current reconstruction. A weighted Mertens capacity energy $A_N$ is introduced and proved to satisfy $A_N/N\to\infty$ and $N/A_N\to0$. The power-growth exponent of the excess provides a critical coordinate. After classical identification, zero transverse escape, local-uniform Mobius control, completed-incidence positivity, and subpower capacity excess are proved equivalent to the Riemann Hypothesis in the displayed normalizations. The equivalence theorem is stated boldly and separately from the still-unproved universal estimate required to settle the conjecture. The static five-class Möbius reporter is proved to be a false closure under prime multiplication; the all-scale predictive quotient is the prime-support address plus one absorbing repeated state.
}
\section{Source freeze and classical identification}
The upstream construction is frozen before classical comparison: binary capacity, transparent arities, independent prime coordinates, squarefree incidence orientation, quota topology, nonbinary contrast fibres, null half-transport, and the orientation operator. Only then are the generated scalar character and sign identified with the classical Dirichlet and Mobius objects in their absolute domain.

This paper studies three exact objects: the finite scalarization defect, a weighted Mertens capacity energy, and equivalent formulations of the Riemann Hypothesis. Equivalence is not itself a proof of the conjecture.

\section{Finite scalarization defect}
The finite uncollapsed response group is invertible before scalarization. Scalar multiplication fails for two exact reasons:
\begin{enumerate}
\item repeated-prime collisions;
\item products crossing the shared capacity boundary.
\end{enumerate}
Let $\Phi_N(a;s)$ be the scalarized response and $\lambda_N(s)$ its neutral generator. Differentiation gives a boundary current $R_N(a;s)$.
\begin{theorem}[Boundary-response reconstruction]
For every finite capacity $N$,
\begin{equation}
\boxed{\Phi_N(a;s)=e^{a\lambda_N(s)}\left[1+\int_0^ae^{-t\lambda_N(s)}R_N(t;s)\,dt\right].}
\end{equation}
Every term in $R_N$ is generated by a repeated-coordinate or over-capacity interaction.
\end{theorem}
Finite zeros therefore satisfy an exact coupled charge-current constraint. The theorem identifies the missing scalar information; it does not estimate the infinite signed current.

\section{Weighted Mertens capacity energy}
Let
\begin{equation}
M(x)=\sum_{n\le x}\mu(n),
\end{equation}
and define
\begin{equation}
A_N=\sum_{q\le N}M\!\left(\left\lfloor\frac Nq\right\rfloor\right)^2,
\qquad \lambda_N=A_N^{-1}.
\end{equation}
Grouping equal values of $\lfloor N/q\rfloor$ gives
\begin{equation}
\frac{A_N}{N}=\sum_{k=1}^NM(k)^2
\frac{\lfloor N/k\rfloor-\lfloor N/(k+1)\rfloor}{N}.
\end{equation}
For fixed $k$, the coefficient tends to $1/[k(k+1)]$.

\begin{theorem}[Divergence of the weighted Mertens energy]
\begin{equation}
\sum_{n\ge1}\frac{M(n)^2}{n(n+1)}=\infty.
\end{equation}
Consequently,
\begin{equation}
\boxed{A_N/N\to\infty,\qquad N\lambda_N\to0.}
\end{equation}
\end{theorem}
\begin{proof}[Proof sketch]
If the weighted series converged, then $g(u)=e^{-u/2}M(e^u)$ would lie in $L^2(0,\infty)$. Its Laplace transform would have uniformly bounded vertical $L^2$ norms in the right half-plane. Summation by parts identifies the transform with $[(z+1/2)\zeta(z+1/2)]^{-1}$ in the initial domain. A critical-line zero gives boundary growth incompatible with the uniform bound. The lower limit for $A_N/N$ then dominates arbitrarily long partial sums of the divergent weighted energy.
\end{proof}

\begin{table}[H]\centering
\caption{Independent finite audit of the capacity energy.}
\begin{tabular}{rrrr}
\toprule
$N$ & $M(N)$ & $A_N$ & $A_N/N$\\\midrule
10 & -1 & 11 & 1.100000\\
100 & 1 & 143 & 1.430000\\
$10^3$ & 2 & 1479 & 1.479000\\
$10^4$ & -23 & 15829 & 1.582900\\
$10^5$ & -48 & 157055 & 1.570550\\
$10^6$ & 212 & 1645995 & 1.645995\\
\bottomrule
\end{tabular}
\end{table}
Finite nonmonotonicity is compatible with the asymptotic theorem.

\section{Capacity recurrence exponent}
Define the critical excess and its power-growth exponent by
\begin{equation}
\mathcal C_N=A_N/N,
\end{equation}
\begin{equation}
\chi_{\rm cap}=\limsup_{N\to\infty}\frac{\log\mathcal C_N}{\log N}.
\end{equation}
The unconditional result is $\mathcal C_N\to\infty$. The critical distinction is subpower growth versus positive-power growth.

\section{Equivalent completion statements}
After the classical identification, the following formulations are equivalent to the Riemann Hypothesis within the frozen normalizations used here:
\begin{enumerate}
\item local-uniform boundedness of the Mobius polynomials $F_N(s)=\sum_{n\le N}\mu(n)n^{-s}$ on every compact subset of $\Re s>1/2$;
\item a completed-incidence positivity inequality on the pole-neutral subspace;
\item absence of positive transverse escape in the normalized causal-tail representation;
\item the critical capacity condition $\chi_{\rm cap}=0$.
\end{enumerate}

\begin{theorem}[Equivalence theorem]
The displayed capacity-completion condition is equivalent to the classical Riemann Hypothesis.
\end{theorem}
\claimstatus{Exact equivalence. The universal capacity estimate required to conclude the condition is not proved here.}

\section{Completed-incidence formulation}
Let $D_L$ be the completed incidence operator and $\Delta(L)$ the common budget in the frozen Weil normalization. On the pole-neutral space $\mathcal H_L^0$, the target inequality is
\begin{equation}
A_L\big|_{\mathcal H_L^0}\ge \Delta(L)I
\qquad\text{for every }L.
\end{equation}
Equivalently, when $\Delta(L)>0$, there is a contraction $T_L$ such that
\begin{equation}
\sqrt{\Delta(L)}\,g=T_LD_Lg,
\qquad \|T_L\|\le1.
\end{equation}
The operator and scale-consistency identities are constructed; the universal lower bound is the open RH-equivalent statement.

\section{Causal-tail formulation}
Let $\Theta_N$ be the normalized adverse tail over degree at most $N$ beyond a sublinear observer boundary. Define
\begin{equation}
\chi_{\rm RH}=\limsup_{N\to\infty}\frac1N\log(1+\Theta_N).
\end{equation}
Positive $\chi_{\rm RH}$ represents exponentially coherent adverse witnesses across increasing degree. The critical regime is the absence of positive exponential escape. This is not identical to the Schatten boundary $\Re s=1/2$ without the explicit intertwiner supplied by the causal-tail construction.

\section{Failed bridges retained}
The following routes are not used as proof:
\begin{itemize}
\item finite zero audits;
\item symmetry of the functional equation alone;
\item random-sign or square-root-cancellation assumptions;
\item a fitted spectrum or random-matrix analogy;
\item the unbounded uniform odd trace on the Hilbert-Schmidt carrier;
\item finite positivity without a uniform bound.
\end{itemize}
They remain valuable countermodels because they identify exactly where a plausible shortcut loses the required scalar orientation.

\section{Static five-class reporter and all-prime predictive closure}
The digit carrier admits the exact terminal-resolved classes
\begin{equation}
\{0\},\quad\{1\},\quad\{2,3,5,7\},\quad\{6\},\quad\{4,8,9\},
\end{equation}
corresponding to zero, unit, M\"obius value $-1$, value $+1$, and value $0$. This is the coarsest static reporter on the ten digits that preserves terminal role and M\"obius value. It is not an autonomous all-scale arithmetic state.

\begin{theorem}[Five-class false closure under prime multiplication]
When multiplication by arbitrary primes is admitted, the five-class reporter is not a continuation congruence. The coarsest predictive quotient for all future M\"obius reports consists of one state for every finite squarefree prime-support set and one absorbing repeated-prime state.
\end{theorem}
\begin{proof}
The states $2$ and $3$ share the prime class. Multiplication by $2$ sends them to $4$ and $6$, whose M\"obius values are $0$ and $+1$. Hence the static class is not future-complete. Represent a squarefree integer by its finite prime support and every nonsquarefree integer by $\dagger$. Multiplication by $p$ adjoins $p$ if absent and enters $\dagger$ if present; $\dagger$ remains absorbing. If two supports differ, multiplying by a prime in their symmetric difference makes one repeated and the other squarefree. Thus distinct supports are future-distinguishable and the quotient is coarsest.
\end{proof}

The source has complete vector and mean-square realizations above $\Re s=1/2$, while fixed-phase scalar evaluation remains unbounded. The exact RH gate is therefore not another finite classification. It is local-uniform boundedness of the special M\"obius diagonal on every compact subset of $\Re s>1/2$.

\section{Conclusion}
The arithmetic programme supplies a source-ordered derivation, an exact finite response reconstruction, an unconditional coercivity theorem, and several exact equivalent forms of RH. It does not yet supply the universal capacity-completion estimate. The live mathematical object is therefore precise: prove subpower critical excess or the equivalent completed-incidence bound without importing the desired cancellation.


\section{Weighted capacity energy in detail}
Let $M(x)=\sum_{n\le x}\mu(n)$ and define
\begin{equation}
A_N=\sum_{q\le N}M\!\left(\left\lfloor\frac Nq\right\rfloor\right)^2,
\qquad \lambda_N=A_N^{-1}.
\end{equation}
Grouping equal quotient values yields
\begin{equation}
\frac{A_N}{N}=
\sum_{k=1}^NM(k)^2
\frac{\lfloor N/k\rfloor-\lfloor N/(k+1)\rfloor}{N}.
\end{equation}
For fixed $k$, the coefficient tends to $1/[k(k+1)]$.

Suppose, for contradiction, that
\begin{equation}
\sum_{n\ge1}\frac{M(n)^2}{n(n+1)}<\infty.
\end{equation}
Then $g(u)=e^{-u/2}M(e^u)$ is square integrable. Its Laplace transform has uniformly bounded vertical $L^2$ norm in the open right half-plane. In the initial domain, summation by parts identifies it with the reciprocal zeta factor divided by the Mellin variable. A zero on the critical line produces boundary blow-up incompatible with that uniform bound. Hence the weighted Mertens energy diverges. Fatou's lemma applied to the grouped expression gives
\begin{equation}
\frac{A_N}{N}\longrightarrow\infty,
\qquad N\lambda_N\longrightarrow0.
\end{equation}
The finite audit is
\begin{center}
\begin{tabular}{rrrr}
\toprule
$N$&$M(N)$&$A_N$&$A_N/N$\\\midrule
10&-1&11&1.100000\\
100&1&143&1.430000\\
$10^3$&2&1479&1.479000\\
$10^4$&-23&15829&1.582900\\
$10^5$&-48&157055&1.570550\\
$10^6$&212&1645995&1.645995\\
\bottomrule
\end{tabular}
\end{center}
Finite nonmonotonicity is compatible with the asymptotic theorem.

\section{Three equivalent critical coordinates}
The classical target can be written through several exact equivalences after the arithmetic identification is frozen.
\begin{enumerate}
\item Local-uniform boundedness of the M\"obius Dirichlet polynomials on compact subsets of $\Re s>1/2$.
\item A completed-incidence positivity or coercivity inequality on every finite scale.
\item Zero power-growth exponent
\begin{equation}
\chi_{\rm cap}=\limsup_{N\to\infty}\frac{\log(A_N/N)}{\log N}=0.
\end{equation}
\end{enumerate}
The unconditional divergence $A_N/N\to\infty$ and the critical condition $\chi_{\rm cap}=0$ are consistent: the latter permits divergent growth slower than every positive power.

\section{Finite scalarization reconstruction}
Before scalarization, the squarefree response group is everywhere invertible. Scalar multiplication fails only when prime coordinates repeat or when a product crosses the shared capacity boundary. Differentiation at the neutral response gives a boundary current $R_N(a;s)$ and exact reconstruction
\begin{equation}
\Phi_N(a;s)=e^{a\lambda_N(s)}
\left[1+\int_0^ae^{-t\lambda_N(s)}R_N(t;s)\,dt\right].
\end{equation}
Every finite scalar zero therefore obeys a coupled charge--current equation whose terms are generated by those two failures. Passing from that exact finite equation to an infinite signed completion is the live analytic problem.

\section{Countermodels and proof boundary}
The source programme tested and retained countermodels to several tempting shortcuts: finite positivity without uniform control; quadratic completion without an odd scalar; exact reflection without a magnitude estimate; sparse bad modes without control of coherent near-resonance; and numerical zero alignment without a source-generated universal inequality. The paper proves equivalence and several auxiliary theorems. It does not claim the open universal bound.

\begin{thebibliography}{99}
\bibitem{titch} E. C. Titchmarsh, revised by D. R. Heath-Brown, The Theory of the Riemann Zeta-Function, Oxford University Press, 1986.
\bibitem{iwakow} H. Iwaniec and E. Kowalski, Analytic Number Theory, American Mathematical Society, 2004.
\bibitem{montgomery} H. L. Montgomery and R. C. Vaughan, Multiplicative Number Theory I: Classical Theory, Cambridge University Press, 2007.
\bibitem{balazard} M. Balazard, E. Saias, and M. Yor, Notes sur la fonction de Riemann, 2, Advances in Mathematics 143 (1999), 284--287.
\bibitem{nyman} B. Nyman, On the One-Dimensional Translation Group and Semi-Group in Certain Function Spaces, thesis, Uppsala, 1950.
\end{thebibliography}
\end{document}
